\taughtsession{Lecture}{Network Layer Security}{2026-04-28}{09:00}{Fahad}{}

Some content in this lecture overlaps significantly with an \textit{IT \& Internetworking Security} lecture from a few weeks ago.

\section{IPsec}
The \textit{Internet Protocol security} framework is a security architecture for IP. There are three key components, some of which each have options for implementations, these are used in combination to build up the framework:
\begin{itemize}
    \item Security Protocol: Either Authentication Header (AH) or Encapsulating Security Payload (ESP)
    \item Key Management: manual or automated (using the Internet Key Exchange)
    \item Cryptographic Algorithms, for authentication and encryption
\end{itemize}

As we can predict, IPsec provides confidentiality, integrity, authentication and anti-replay mechanisms. \textit{Confidentiality} is provided through encryption of data meaning only the sender and receiver can read the data. \textit{Integrity} is provided by calculating and verifying hashes to enable the receiver to check the packet was not modified in transit. \textit{Authentication} is enabled through ensuring that the sender and receiver authenticate each other to make sure that they are communicating with the intended entity. \textit{Anti-replay} prevents duplicate packet transmission, and is supported in IPsec through assigning each packet a sequence number. 

The cryptographic algorithms, as mentioned above, are not explicitly stated in the IPsec framework. Rather the framework leaves space for a modular system of protocols to be inserted depending on the requirements of the network. However, the selected algorithms must come from a list of approved algorithms such that they work with AH and ESP; similarly for IKEv2.

Due to each component of IPsec being specified under separate intertwined specifications, a single component can be updated without the entire specification having to be re-written.

IPsec creates a boundary between unprotected and protected interfaces for a host or a network. For example on a edge router between a LAN and the Internet.  Then when traffic wants to cross this boundary, one of three things happens:
\begin{itemize}
    \item traffic passes unimpeded
    \item traffic is afforded security services via AH or ESP and sent on its merry way
    \item traffic is discarded
\end{itemize}

IPsec is used to protect traffic in numerous different scenarios:
\begin{itemize}
    \item between a pair of hosts
    \item between a pair of security gateways - as in a site-to-site connection over the internet
    \item between a security gateway and a host - as in a teleworker connecting to the head office
\end{itemize}


\subsection{Security Protocols}
\textit{Authentication Header} (AH) provides data integrity, authentication of IP packets, and protection against replay attacks. The first two are achieved through the use of a Message Authentication Code (MAC), i.e. HMAC-SHA1-96. IP packets expand when AH added with item such as:
\begin{itemize}
    \item next header - type of the header following this header
    \item payload length - length of AH
    \item Security Parameter Index (SPI) - identifies a SA
    \item sequence number
    \item authentication data
\end{itemize}

\textit{Encapsulating Security Protocol} (ESP) provides the functionality of AH, plus confidentiality. When applied to an IP packet, ESP adds an ESP header, encrypts the payload, and adds an ESP trailer. An ESP packet contains:
\begin{itemize}
    \item security parameter index (SPI)
    \item sequence number
    \item payload data (encrypted)
    \item padding - to achieve data length that is a multiple of 32 bits (encrypted)
    \item padding length (encrypted)
    \item next header (encrypted)
    \item (optional) authentication header
\end{itemize}

\subsection{Modes of Operation}
There are two modes of operation: \textit{transport mode} and \textit{tunnel mode}. A host implementation of IPsec must support both transport and tunnel modes; while a security gateway implementation must support tunnel mode and may support transport mode.

\textit{Transport mode} is the least secure of the two. Only the payload of the IP packet is always protected; with the header not encrypted in ESP, but some components of it encrypted in AH. The data is protected from source to destination, but the header information is transmitted as cleartext. This is used predominantly between hosts. 

\textit{Tunnel Mode} is the more secure of the options. The entire IP packet is encrypted (this meaning both the IP header and payload), and this becomes the payload of a new IP packet. The new (outer) IP packet may contain different source and destination addresses, obfuscating the original source and destination data. This provides data flow confidentiality to some extent and can be used in gateway-gateway, host-host, or host-gateway paradigms. 

\subsection{Methods of IPsec}

There are four combinations of security protocol and operating mode. These all come in context of the original IP packet, as below.
\begin{figure}[H]
    \centering
    \begin{tikzpicture}
    \node[name=data, rectangle split, 
         rectangle split horizontal, 
         rectangle split parts=3, 
         rectangle split draw splits=true, 
         draw,
         text width=2.5cm,
         align=center,
         font=\strut
         ] (main) 
        {IP Header \nodepart{two} TCP Header 
              \nodepart{three} Data};

    \end{tikzpicture}
    \caption{Original IP Packet}
\end{figure}

\subsubsection{Authentication Header: Transport Mode}
In transport mode, while using AH, the HMAC is computed over:
\begin{itemize}
    \item immutable header fields (the fields that do not change in transit) which include source address, IP header length
    \item the AH (except authentication data field)
    \item IP data
\end{itemize}
The values for the excluded fields are set to zero during HMAC computation.

\begin{figure}[H]
    \centering
    \begin{tikzpicture}
    \node[name=data, rectangle split, 
         rectangle split horizontal, 
         rectangle split parts=4, 
         rectangle split draw splits=true, 
         draw,
         text width=2.5cm,
         align=center,
         font=\strut
         ] (main) 
        {IP Header \nodepart{two} AH Header 
              \nodepart{three} TCP Header
              \nodepart{four} Data};

        \draw[<->, >=stealth, thick] 
            ([yshift=-2mm]main.south west) -- 
            ([yshift=-2mm]main.south east) 
            node[midway, below] {Authenticated};

    \end{tikzpicture}
    \caption{AH Transport Mode}
\end{figure}

\subsubsection{Authentication Header: Tunnel Mode}
In tunnel mode while using AH, the HMAC is computed over:
\begin{itemize}
    \item immutable header fields (fields that do not change in transit) which include source address, IP header length
    \item the AH (except authentication data fields)
    \item IP data
\end{itemize}
The values for excluded fields are set to zero during HMAC computation.

\begin{figure}[H]
    \centering
    \begin{tikzpicture}
    \node[name=data, rectangle split, 
         rectangle split horizontal, 
         rectangle split parts=5, 
         rectangle split draw splits=true, 
         draw,
         text width=2.5cm,
         align=center,
         font=\strut
         ] (main) 
        {New IP Header \nodepart{two} AH Header
              \nodepart{three} IP Header 
              \nodepart{four} TCP Header
              \nodepart{five} Data};

        \draw[<->, >=stealth, thick] 
            ([yshift=-2mm]main.south west) -- 
            ([yshift=-2mm]main.south east) 
            node[midway, below] {Authenticated};

    \end{tikzpicture}
    \caption{AH Tunnel Mode}
\end{figure}

\subsubsection{Encapsulating Security Payload: Transport Mode}
Authentication is only provided over ESP header and data fields. Only the data is encrypted, meaning the source and destination are clear.

\begin{figure}[H]
    \centering
    \begin{tikzpicture}
    \node[name=data, rectangle split, 
         rectangle split horizontal, 
         rectangle split parts=6, 
         rectangle split draw splits=true, 
         draw,
         text width=2.5cm,
         align=center,
         font=\strut
         ] (main) 
        {IP Header \nodepart{two} ESP Header
              \nodepart{three} TCP Header 
              \nodepart{four} Data
              \nodepart{five} ESP Trailer
              \nodepart{six} ESP Auth};

        \draw[<->, >=stealth, thick] 
            ([yshift=-2mm]main.two split south) -- 
            ([yshift=-2mm]main.five split south) 
            node[midway, below] {Encrypted};

        \draw[<->, >=stealth, thick] 
            ([yshift=-8mm]main.one split south) -- 
            ([yshift=-8mm]main.five split south) 
            node[midway, below] {Authenticated};

    \end{tikzpicture}
    \caption{ESP Transport Mode}
\end{figure}

\subsubsection{Encapsulting Security Payload: Tunnel Mode}
Authentication is only provided over ESP Header and Data, not over IP header fields. The entire payload is encrypted, meaning that the real source and destination addresses are hidden. 

\begin{figure}[H]
    \centering
    \begin{tikzpicture}
    \node[name=data, rectangle split, 
         rectangle split horizontal, 
         rectangle split parts=7, 
         rectangle split draw splits=true, 
         draw,
         text width=2cm,
         align=center,
         font=\strut
         ] (main) 
        {New IP Header \nodepart{two} ESP Header
              \nodepart{three} IP Header
              \nodepart{four} TCP Header 
              \nodepart{five} Data
              \nodepart{six} ESP Trailer
              \nodepart{seven} ESP Auth};

        \draw[<->, >=stealth, thick] 
            ([yshift=-2mm]main.two split south) -- 
            ([yshift=-2mm]main.six split south) 
            node[midway, below] {Encrypted};

        \draw[<->, >=stealth, thick] 
            ([yshift=-8mm]main.one split south) -- 
            ([yshift=-8mm]main.six split south) 
            node[midway, below] {Authenticated};

    \end{tikzpicture}
    \caption{ESP Tunnel Mode}
\end{figure}

\section{Virtual Private Networks}
A \textit{Virtual Private Network} (VPN) is a method of securely communicating data over an unsecured (or public) network. They are used in enterprise to extend connections to locations beyond the corporate trusted network. For example a teleworker uses a VPN to connect to the corporate network to access the file server, meaning the file server then doesn't have to be directly exposed to the internet. 

VPNs must accomplish three things:
\begin{itemize}
    \item Encapsulation of incoming and outgoing data
    \item Encryption of incoming and outgoing data (for secure VPNs)
    \item Authentication of remote computers and remote users
\end{itemize}

Similar to IPsec, a VPN can operate in either transport or tunnel mode.

\subsection{Transport Mode VPN}
Transport mode VPNs encrypt the data within the IP packet, but not the header of the packet; meaning that the origin and destination addresses are visible should the packet be inspected. 

These VPNs allow for users to connect to a remote network directly from their workstation, with the network acting as though their workstation is directly connected to the remote network. This is the popular paradigm for remote access workers to use when connecting to their office network over the Internet, by connecting to a VPN server on the perimeter of the corporate network.

The VPN server acts as intermediate client and decrypts traffic coming from the client / encrypts traffic going to the client. The destination machine on the corporate network receives encrypted data and decrypts it.

\subsection{Tunnel Mode VPN}
Tunnel mode VPNs operate between two perimeter tunnel servers. These act as encryption points, encrypting all traffic that will traverse the unsecure network. On the receiving network, all traffic is decrypted by the perimeter tunnel server and released into the network in a decrypted state. This means that the end stations receive decrypted data. 

The primary benefit of this model is that an intercepted packet reveals nothing about the true destination system. 

\section{Internet Key Exchange}
The \textit{Internet Key Exchange} (IKE) is a protocol used within IPsec to set up authenticated and encrypted communication channels between two devices which is used for VPN traffic. IKE is used to agree the parameters used in an IPsec communication - AH vs ESP, establish symmetric key, and cryptographic parameters used. IKE is the key management layer of IPsec. There are two versions: IKEv1 and IKEv2.

\subsection{IKEv1}
IKEv1 has two phases - the first creating the \textit{IKE SA} (IKE Phase 1) tunnel between the two connected devices. The second creates an inner \textit{IPsec SA} (IKE Phase 2) tunnel within the IKE Phase 1 tunnel. The data travels through this inner tunnel. 

Phase 1 has two methods of operation: main mode and aggressive mode.

\begin{extlink}
There are diagrams showing what is written out below available in the slides available on Moodle.
\end{extlink}

\subsubsection{Phase 1: Main Mode}
Main mode of phase 1 establishes a secure, authenticated and encrypted channel between peers using a six-message exchange:
\begin{enumerate}
    \item A $\rightarrow$ B: Send the local IKE proposal
    \item B $\rightarrow$ A: Confirm the IKE proposal
    \item A $\rightarrow$ B: Send key-related information
    \item B $\rightarrow$ A: Send key-related information
    \item A $\rightarrow$ B: Send identity information
    \item B $\rightarrow$ A: Send identity information
\end{enumerate}

The first two messages are used to negotiate the parameters for establishing the IKE SA. A sends B the proposal of how it would like to communicate, offering a number of options, to which B responds with the chosen SA. 

The third and fourth messages are used to exchange key related information and generate the key. 

The final pair of messages, fifth and sixth, are used to exchange identity information to authenticate the peer.

See below `\textit{Security Associations}' for more information on the HAGLE mnemonic and the process it outlines.

\subsubsection{Phase 1: Aggressive Mode}
Aggressive mode of phase 1 accelerates the same activities that main mode uses. This is used to establish tunnels faster, as it cuts down the number of required packets from 6 to 3. Commonly seen in the scenario where one peer has a dynamic IP address or where multiple clients connect to a single VPN gateway. 

\begin{enumerate}
    \item A $\rightarrow$ B: Send the local IKE proposal, key related information and identity information
    \item B $\rightarrow$ A: Search for the matching IKE proposal. Sends the matching IKE proposal, key-related information, identity information and local authentication information
    \item A $\rightarrow$ B: Responds with the local authentication information to implement authentication
\end{enumerate}

Aggressive mode is less secure than main mode as it transmits identification information in plain text without encryption.

\subsubsection{Phase 2}
IKEv1 Phase 2, known as quick mode, is used to negotiate the security parameters of the inner IPsec tunnel which is used to protect the user data passing through the tunnels. It is a three-packet configuration:
\begin{enumerate}
    \item A $\rightarrow$ B: Sends the IPsec parameters adn identity authentication information
    \item B $\rightarrow$ A: Sends the matching IPsec parameters and identity information and generates keys
    \item A $\rightarrow$ B: Sends the confirmation information and complete phase 2 negotiation
\end{enumerate}

\subsection{Security Associations}
During phase 1 of establishing a tunnel between two peers, there are 5 steps which must be undertaken. The steps can be remembered with the HAGLE mnemonic.
\begin{itemize}
    \item[H] Hash Algorithm - HMAC or Hashed Message Authentication Mode
    \item[A] Authentication Method - how peers prove their identity (more below)
    \item[G] Diffie-Hellman Group - determines the strength of the keys used in the key exchange process
    \item[L] Lifetime - Defines how long the Phase 1 SA remains active before it must be re-negotiated
    \item[E] Encryption Algorithm - Specifies the algorithm used to encrypt the data 
\end{itemize}

Both sides must agree on the parameters to use and they must match, else the IPsec tunnel will fail to establish. This only applies to the Phase 1 connection - creating the management tunnel, not the phase 2 data tunnel. 

\subsubsection{Diffie-Hellman Key Exchange}
The Diffie-Hellman (DH) key exchange uses a series of mathematical operations to identity a shared secret key. It uses the modular exponentiation (power mod), represented as MODP.

\begin{enumerate}
    \item A and B agree modulus $p$ and base $g$
    \item A chooses secret integer $a$, then sends B: $X=g^a \mod p$
    \item B chooses secret integer $b$, then sends A: $Y=g^b \mod p$
    \item A computes $Z=Y^a \mod p$
    \item B computes $Z=X^b \mod p$
    \item A and B now share the secret $Z$
\end{enumerate}

This works because $X^b \mod p = g^{ab} \mod p = g^{ba} \mod p = Y^a \mod p$

\begin{example}{DH Key Exchange}
A and B agree modulus $p=23$ and base $g=5$

A chooses secret integer $a=6$ then sends B:
\[X = g^a \mod p = 5^6 \mod 23 = 8\]

B chooses secret integer $b=15$, then sends A:
\[Y=g^b \mod p = 5^{15} \mod 23 = 19\]

A computes:
\[Z=Y^a \mod p = 19^6 \mod 23 = 2\]

B computes:
\[Z=X^b \mod p = 8^{15} \mod 23 = 2\]

A and b now share a secret (the number 2).
\end{example}

\subsubsection{Authentication}

\begin{todo}
Come back to this, look at `diagrams' on Moodle slides and videos to make sense of it.
\end{todo}

\subsection{IKEv2}
IKEv2 is a revised version of IKEv1. It revises 12 things:
\begin{enumerate}
    \item Define entire protocol in 1 document, replacing 2 RFCs; incorporate changes to support NAT Traversal, Extensible Authentication and remote address acquisition
    \item Simplify IKE by replacing the eight different initial exchanges with single four-message exchange
    \item Remove the domain of interpretation, situation and labelled domain identifier fields, and the commit and authentication bits
    \item Decrease IKE's latency in the common case by making the initial exchange 4 messages, and allowing piggybacking setup of Child SA on that exchange
    \item Replace the cryptographic syntax for protecting IKE messages with one based closely on ESP to simplify implementation and security analysis
    \item Reduce the number of possible error states by making the protocol reliable (i.e. all messages are acknowledges) and sequenced; has allowed for shortening creation of a child SA from 3 messages to 2
    \item Increase robustness by allowing the responder to do not significant processing until it receives a message providing that the initiator can received messages at its claimed IP address
    \item Fix cryptographic weaknesses such as the problem with symmetries in hashes used for authentication
    \item To specify Traffic Selectors in their own payload type rather than overloading ID payloads; making more flexible the required Traffic Selectors
    \item Specify required behaviors under certain error conditions or when data that is not understood is received in order to make it easier to make future revisions in a way that does not break backwards compatibility
    \item Simplify and clarify how shared state is maintained in the presence of network failures and DoS attacks
    \item Maintain existing syntax and magic numbers to the extent possible to make it likely that implementations of IKEv1 can be enhances to support IKEv2 with minimum effort.
\end{enumerate}

The IKEv2 tunnel is a single layered tunnel and is established through a single four-message exchange:
\begin{enumerate}
    \item A $\rightarrow$ B: IKE SA initial exchange
    \item B $\rightarrow$ A: IKE SA initial exchange
    \item A $\rightarrow$ B: Authentication exchange
    \item B $\rightarrow$ A: Authentication exchange
\end{enumerate}

The firs two messages are used for IKE SA establishment and key generation. The final two messages are usd to identify authentication and the establishment of rhe first pair of IPsec SAs. 